\par\noindent\SourceLeaf{65}%
Dividiert man diesen Ausdruck durch $0$, so bleibt
\[
\frac{1}{2^{2h}}(-1)^{2h-1}\frac{1.2.\,.\,.\,.(2h-1)}{1.2.\,.\,.\,. 2h}\,{}^{2h}\mathfrak{B}^{(h)}\,n^{2h} = (-1)^{2h-1}\frac{1}{2^{2h}} . \frac{1}{2h}\,{}^{2h}\mathfrak{B}^{(h)}\,n^{2h}
\]
\[
= \frac{1}{2^{2h}}(-1)^{2h-1} . \frac{(2h-1)(2h-2).\,.\,.(h+1)}{1.2.\,.\,. h}\,n^{2h}
\]
\[
= \frac{1}{2^{h}}(-1)^{2h-1}\frac{1.3.5.\,.\,.(2h-1)}{1.2.3.\,.\,.\,. h.}\,\frac{n^{2h}}{2h}
\]
oder, da $(-1)^{2h-1}$ für jeden Wert von $h$ negativ ist,
\[
= -\frac{1}{2^{h}}\frac{1.3.5.\,.\,.(2h-1)}{1.2.3.\,.\,. h}\,\frac{n^{2h}}{2h}.
\]
Differenziert man diesen Ausdruck nach $n$ und lässt $\mathrm{d}n$ weg, so ergibt sich
\[
-\frac{1}{2^{h}}\frac{1.3.5.\,.(2h-1)}{1.2.3.\,.\,. h}\,n^{2h-1},
\]
Dies ist offenbar das Glied der Größe
\[
\frac{-(1-nn)^{-\frac{1}{2}}}{n}
\]
mit Index $h^{\mathrm{tus}}$. Da also das Anfangsglied
$\dfrac{-1}{n}$ jener Reihe hier fehlt, ist sie nach Abzug dieses Glieds unserer Formel gleich. Schreibt man demnach den
Anfangskoeffizienten von $\log(1+n\cos\varphi)$ in der Form $=C$, so erhält man
\[
\frac{n\,{}^{*}\!\mathrm{d}C}{\mathrm{d}n} = -(1-nn)^{-\frac{1}{2}}+1.
\]
Schreibt man die Gleichung in der Form
\[
{}^{n}\mathrm{d}C = \frac{\mathrm{d}n}{n}\bigl(1-(1-nn)^{-\frac{1}{2}}\bigr),
\]
so ergibt die Integration
\[
C = \log n - \log\frac{1-(1-nn)^{\frac{1}{2}}}{n} + \mathrm{const};
\]
Da der Ausdruck für $n=0$ verschwinden muss, ist die Konstante $=\log\tfrac{1}{2}$. Daher
\[
C = \log\left(\frac{nn}{(1-\sqrt{1-nn})2}\right),
\]
oder, wenn man es vorzieht,
\[
C = -\log\left(\frac{2(1-\sqrt{1-nn})}{nn}\right).
\]

